\documentclass[10pt,a4paper]{article}

\usepackage[top=1.7cm,bottom=1.8cm,left=1.8cm,right=1.8cm]{geometry}
\usepackage[T1]{fontenc}
\usepackage[utf8]{inputenc}
\usepackage{graphicx}
\usepackage{float}
\usepackage{xcolor}
\usepackage{caption}
\usepackage{hyperref}
\usepackage{titlesec}

\definecolor{ink}{HTML}{252525}
\definecolor{muted}{HTML}{66645F}
\definecolor{accent}{HTML}{2878D0}
\definecolor{soft}{HTML}{F3F5F7}
\definecolor{quotebg}{HTML}{FBF7EC}

\hypersetup{colorlinks=true,urlcolor=accent,linkcolor=accent}
\setlength{\parindent}{0pt}
\setlength{\parskip}{0.5em}
\captionsetup{font=footnotesize,labelfont=bf,skip=5pt}
\titleformat{\section}{\Large\bfseries\color{ink}}{\thesection}{0.7em}{}
\titleformat{\subsection}{\large\bfseries\color{ink}}{\thesubsection}{0.6em}{}

\newcommand{\paperquote}[1]{%
  \par\vspace{0.15em}%
  \noindent\colorbox{quotebg}{\parbox{0.97\textwidth}{\footnotesize\itshape
  From the atlas paper: ``#1''}}\par\vspace{0.15em}}

\newcommand{\netdef}{A field of view (FOV) is \emph{eligible} for a rule only when it holds
at least 20 cells of both named types. Among eligible FOVs the rule is attraction, avoidance,
or no rule, and the net score is
$100(N_{\rm attraction}-N_{\rm avoidance})/N_{\rm eligible}$.}

% The four-panel figure every single rule gets. Written out once.
\newcommand{\summaryvalues}{Four views of one rule, all on the same eligible FOVs.
\textbf{How often it fires:} the share of that stage's eligible FOVs scoring attraction and
the share scoring avoidance, as two separate lines, with the eligible count below.
\textbf{What it is made of:} one dot per eligible FOV giving that FOV's share of cells
belonging to each named type, with the stage mean as a diamond; the denominator is every cell
in the FOV, not just the two types. \textbf{How strong:} one dot per rule-bearing FOV and the
stage median, where Lift is observed co-occurrence divided by what independence predicts, so
1 is neutral, above 1 supports attraction and below 1 supports avoidance; the count below is
rule-bearing / eligible. \textbf{Every FOV:} all FOVs of the stage split into attraction,
no rule, avoidance and too few cells to judge, so composition loss stays visible instead of
being scored as ``no rule''; the net score sits above each bar. \netdef}

\newcommand{\fovvalues}{The top row shows every cell in the FOV; the bottom row greys out all
cells except the two named types. Within each stage the FOV is picked automatically: take the
more frequent rule direction among eligible rule-bearing FOVs, then the occurrence whose Lift
is closest to that stage's median. The state, Lift and mining FDR above the bottom map belong
to that one FOV, not to the between-stage test. These maps illustrate; the aggregate figure
before them is the evidence.}

\newcommand{\heatvalues}{Rows are rules, ordered by how often they occur. Each cell gives the
count and the percentage of units in that group holding the rule. The bars on the left are the
mean share of cells of the antecedent (left) and consequent (right) type, over the FOVs where
the rule fires. No eligibility filter: every unit is counted.}

\newcommand{\dumbbellvalues}{Each row is one rule; the two dots are its net score in the two
groups and the bar is the gap in percentage points. Hollow dots and grey text mark rules that
did not pass FDR 0.05. Rows are ordered by gap size.}

\newcommand{\withcellsvalues}{Left: the net score of each rule that moves steadily across the
three stages, with its change per stage beside the name. Right: the mean share of the cell
types those same rules name, with their own change per stage. Both panels are in percentage
points, so the two can be compared directly. A dashed line moved back at one step instead of
moving throughout. No eligibility filter.}

\newcommand{\trendvalues}{Each line is one rule's net score among eligible FOVs across the
three stages, with its change per stage beside the name. Only rules passing FDR 0.05 in the
ordered Control--Mild--Severe test are drawn, so an empty panel means nothing passed.}

\newcommand{\countvalues}{Every dot is one FOV's raw count of the named cell type, with the
20-cell eligibility threshold as a dashed line and the number of FOVs reaching it printed
above each stage. This panel decides what the rule panels can measure at all.}

\newcommand{\statelift}{Each dot is one rule-bearing FOV, restricted to occurrences of a
single direction, with the stage median as a diamond. For avoidance, lower Lift means
stronger separation. The count below is rule-bearing / eligible FOVs.}

\begin{document}

\begin{center}
  {\LARGE\bfseries Differential pair rules in gastrointestinal aGVHD}\\[0.35em]
  {\large What the method resolves, what agrees with the atlas, and what is new}
\end{center}
\vspace{0.2cm}
\hrule
\vspace{0.4cm}

\colorbox{soft}{\parbox{0.97\textwidth}{
\textbf{What is measured.} For every ordered pair of cell types and every FOV, the mining step
asks whether the two types sit closer together than chance (attraction), further apart than
chance (avoidance), or neither. \netdef{} Direction matters: A $\rightarrow$ B and
B $\rightarrow$ A are separate rules with their own counts and their own tests. The run
behind this document holds 525 directional pair rules over 267 FOVs.

\textbf{Which severity score.} Stages are \textbf{clinical} unless a figure says otherwise,
because that is the score the atlas paper uses for its own immune analyses. The pathological
score is shown wherever the two disagree, and \hyperlink{scores}{they disagree a great deal} --
which turns out to be one of the findings rather than a nuisance.

\hypertarget{test-guide}{\textbf{How to read the tests.}}
``Ordered FDR'' asks whether the net score moves steadily from Control to Mild to Severe.
C--M, M--S and C--S compare the two named stages. FDR is corrected across all rules in that
screen; below 0.05 is called significant. Rule names are never chosen by hand before a screen.

\textbf{Which trends count.} A result is reported here only when Control and Severe differ.
A rule that jumps at Mild and returns to its Control value by Severe is not treated as a
finding, however small its Mild--Severe FDR; those stay in the notebooks.

\textbf{What this cannot say.} These are spatial associations, not mechanisms, between
classifier labels. Many endpoint groups are small. Every figure prints its own denominators,
and those should be read before its FDR.
}}

% =============================================================================
\section{Does the method see real spatial structure?}

One question comes before any biology: does a pair rule measure how cells are \emph{arranged},
or is it a slow way of counting how many cells there are? Nothing in this part depends on a
rule name chosen in advance.

\subsection*{The screen returns a stable, structured ranking}

A few pair rules occur in most fields of view, the rest fall away smoothly, and the stage
columns already separate. Paneth $\rightarrow$ Epithelial drops from 21\% of Control FOVs to
8\% of clinically Severe ones; Endocrine $\rightarrow$ Epithelial climbs from 4\% to 21\%;
CD4T $\rightarrow$ Plasma falls from 38\% to 3\%.

\begin{figure}[H]
  \centering
  \includegraphics[width=0.97\textwidth]{summary_downloads/overview_prevalence_clinical_FOV_bothorgans_prevalence.pdf}
  \caption{\heatvalues}
  \label{fig:heat}
\end{figure}

Counting each patient once instead of each FOV leaves the ordering and the stage pattern
unchanged -- Paneth $\rightarrow$ Epithelial 57\% to 16\% of patients, Endocrine
$\rightarrow$ Epithelial 14\% to 44\%, CD4T $\rightarrow$ Plasma 57\% to 8\% -- so none of
this is an artefact of how many fields each patient contributed.

\begin{figure}[H]
  \centering
  \includegraphics[width=0.97\textwidth]{summary_downloads/overview_prevalence_clinical_PatientID_bothorgans_prevalence.pdf}
  \caption{Values follow Figure~\ref{fig:heat}; one patient is now one unit, so a rule counts
  once however many of that patient's fields hold it.}
\end{figure}

\subsection*{Organization moves much more than composition}

This is the central control. The next two figures use the same FOVs, the same organ and the
same two stages, measured first as rules and then as plain cell counts. The rules move by 37
to 67 percentage points; the cell counts behind them move by at most 9.

\begin{figure}[H]
  \centering
  \includegraphics[width=0.90\textwidth]{summary_downloads/Colon_clinical_Control_vs_Severe_allfovs_dumbbell.pdf}
  \caption{\dumbbellvalues{} Colon, Control against clinically Severe, no eligibility filter.
  See the \protect\hyperlink{test-guide}{test guide above}.}
  \label{fig:dumbbell}
\end{figure}

Measured as plain cell abundance instead, the same fields barely move.

\begin{figure}[H]
  \centering
  \includegraphics[width=0.90\textwidth]{summary_downloads/Colon_clinical_Control_vs_Severe_allfovs_cells_dumbbell.pdf}
  \caption{The same two groups of FOVs and the same layout as
  Figure~\ref{fig:dumbbell}, but each row is now a cell type and the value is its share of all
  cells in the FOV rather than a rule's net score.}
  \label{fig:cells}
\end{figure}

The same contrast holds across all three stages at once. Every colonic rule that climbs is a
goblet or muscle avoidance relaxing by 18 to 31 points per stage, while goblet cells fall 5
points and nothing else moves more than 3.

\begin{figure}[H]
  \centering
  \includegraphics[width=0.97\textwidth]{summary_downloads/Colon_clinical_climbing_with_cells.pdf}
  \caption{\withcellsvalues{} Colon, clinical score. Twenty-two rules climb and the eight
  largest are drawn.}
  \label{fig:withcells}
\end{figure}

The rules that fade tell a sharper story: all eleven end at plasma or CD4T, and all eleven
collapse to zero by Mild and stay there. That is a step at the onset of disease, not a
gradient with severity.

\begin{figure}[H]
  \centering
  \includegraphics[width=0.97\textwidth]{summary_downloads/Colon_clinical_fading_with_cells.pdf}
  \caption{Values follow Figure~\ref{fig:withcells}, for the fading direction.}
\end{figure}

In the duodenum only two rules fade, and one of them moves against its own cells: fibroblast
and muscle both grow, by 3 and 5 points, while Fibroblast $\rightarrow$ Muscle falls by 20
points per stage.

\begin{figure}[H]
  \centering
  \includegraphics[width=0.97\textwidth]{summary_downloads/Duodenum_clinical_fading_with_cells.pdf}
  \caption{Values follow Figure~\ref{fig:withcells}; Duodenum.}
\end{figure}

\subsection*{Losing the cells is not the same as losing the arrangement}

A rule can vanish for two very different reasons: the arrangement broke down, or the cells that
made it are no longer there to count. Every rule therefore gets its own eligibility test, and
the bottom panel keeps the unjudgeable fields visible as a separate grey block. Duodenal
Epithelial $\rightarrow$ Goblet is the clean case: among fields that can still be judged the
rule is as common in Severe disease as in Control (97\%, 88\%, 91\%) and just as strong (Lift
about 1.8 throughout), while the number of judgeable fields falls from 30 of 31 to 23 of 41.

\begin{figure}[H]
  \centering
  \includegraphics[width=\textwidth]{summary_downloads/Duodenum_clinical_Epithelial_to_Goblet_summary.pdf}
  \caption{\summaryvalues{} See the \protect\hyperlink{test-guide}{test guide above}.}
  \label{fig:summary}
\end{figure}

That distinction does real work. Without the eligibility filter, 22 colonic rules climb and 11
fade with clinical severity (Figures~\ref{fig:withcells} and following). With it, \emph{none}
of them passes -- both panels below are empty. Colonic goblet cells fall from 22\% to 13\% of
the tissue, so many fields stop being judgeable, and a rule going from avoidance to unjudgeable
looks exactly like avoidance relaxing.

\begin{figure}[H]
  \centering
  \includegraphics[width=0.97\textwidth]{summary_downloads/Colon_clinical_eligible_trend.pdf}
  \caption{\trendvalues{} Colon, clinical score, eligibility applied.}
  \label{fig:trend}
\end{figure}

The duodenum keeps exactly two rules under the same test. Those two carry most of Part 3.

\begin{figure}[H]
  \centering
  \includegraphics[width=0.97\textwidth]{summary_downloads/Duodenum_clinical_eligible_trend.pdf}
  \caption{Values follow Figure~\ref{fig:trend}; Duodenum.}
  \label{fig:trend-duo}
\end{figure}

\subsection*{The method recovers a boundary already known to be real}

Colon and duodenum are different tissues, and the atlas treats them separately for that reason.
A method that measures organization should rediscover the split without being told. It does,
and all six differences that survive correction are anchored on goblet cells -- the expected
answer, since goblet cells are a fifth of all colonic cells and a twenty-fifth of duodenal ones.

\paperquote{Because of the systematic changes between epithelial and immune composition and
organization in the duodenum and the colon, we proceeded to analyze the duodenal and colonic
aHSCT samples separately, comparing each with their relevant controls.}

\begin{figure}[H]
  \centering
  \includegraphics[width=0.90\textwidth]{summary_downloads/organs_eligible_dumbbell.pdf}
  \caption{Values follow Figure~\ref{fig:dumbbell}; Colon against Duodenum with all stages
  pooled and the eligibility filter applied.}
\end{figure}

\subsection*{Rules are structured across patients}

Summarising each patient by the share of their FOVs holding each rule, and then clustering,
groups patients into recognisable profiles rather than scattering them: one block carries the
stromal and macrophage rules, another the plasma and CD4T rules, and Epithelial $\rightarrow$
Goblet is close to universal.

\begin{figure}[H]
  \centering
  \includegraphics[width=0.97\textwidth]{summary_downloads/overview_fingerprint_patient_clinical_bothorgans.pdf}
  \caption{Rows are rules, columns are patients ordered by how similar their rule profiles are.
  Colour is the share of that patient's FOVs holding the rule. The two strips on top give each
  patient's pathological and clinical stage; the bars on the right are the rule's mean share
  across patients. No eligibility filter.}
\end{figure}

% =============================================================================
\section{What agrees with the atlas}

Where the atlas paper reports a finding, does the rule layer recover it independently? Each
result below is paired with the sentence it corresponds to.

\subsection*{Paneth cells disappear from the duodenal crypt}

Paneth cells sit at the base of the crypt and release antimicrobial peptides. Losing them is
one of the few established tissue markers of gut aGVHD -- the paper uses Paneth counts as the
benchmark its own predictor has to beat. The count panel comes first because it decides what
anything else can measure: the fields holding enough Paneth cells to judge fall from 13 of 31
to 13 of 41.

\paperquote{In accordance with previous reports, we found a reduction in Paneth cells in a
subset of patients but not in all.}

\begin{figure}[H]
  \centering
  \includegraphics[width=0.90\textwidth]{summary_downloads/Duodenum_clinical_Paneth_to_Epithelial_count_Paneth.pdf}
  \caption{\countvalues{} Duodenum, clinical stages.}
  \label{fig:count}
\end{figure}

Here the two severity scores part company, and the difference matters. Under pathological
staging only 6 of 36 Severe fields remain judgeable and the rule's net score falls from
$+77\%$ to $+33\%$; under clinical staging 13 of 41 remain and the net barely moves. Paneth
loss is a histological finding, and the pathological score is the one read off the same slide.

\begin{figure}[H]
  \centering
  \includegraphics[width=\textwidth]{summary_downloads/Duodenum_pathological_Paneth_to_Epithelial_summary.pdf}
  \caption{Values follow Figure~\ref{fig:summary}; pathological stages.}
  \label{fig:paneth-path}
\end{figure}

Under clinical staging the collapse is much milder, and the rising Lift rests on 7
fields rather than 2.

\begin{figure}[H]
  \centering
  \includegraphics[width=\textwidth]{summary_downloads/Duodenum_clinical_Paneth_to_Epithelial_summary.pdf}
  \caption{Values follow Figure~\ref{fig:summary}; the same rule under clinical stages.}
\end{figure}

What is interesting is what happens where Paneth cells survive: their association with
surrounding enterocytes does not weaken, it tightens. The reverse rule puts a corrected test
behind that -- Epithelial $\rightarrow$ Paneth climbs from a net of $0\%$ to $+50\%$, ordered
FDR 0.029 -- though its Severe endpoint is 3 fields out of 6.

\begin{figure}[H]
  \centering
  \includegraphics[width=\textwidth]{summary_downloads/Duodenum_pathological_Epithelial_to_Paneth_summary.pdf}
  \caption{Values follow Figure~\ref{fig:summary}; the reverse rule is counted and tested
  independently of the forward one, under pathological stages.}
\end{figure}

\subsection*{Endocrine cells accumulate where Paneth cells were}

The paper's companion finding is that endocrine cells build up in the crypts, and it raises the
possibility that stem cells misdifferentiate into endocrine cells instead of Paneth cells. The
rule layer shows this as a threshold effect that is hard to argue with: not one control field
in the whole duodenal cohort holds 20 endocrine cells, against 29 of 98 Mild and 8 of 41
Severe.

\paperquote{Whereas in the normal duodenum, the endocrine cells were evenly spread along the
villi, patients undergoing aHSCT showed increased density primarily near the muscularis
mucosa.}

\begin{figure}[H]
  \centering
  \includegraphics[width=0.90\textwidth]{summary_downloads/Duodenum_clinical_Endocrine_to_Epithelial_count_Endocrine.pdf}
  \caption{Values follow Figure~\ref{fig:count}, for endocrine cells.}
\end{figure}

Where the rule becomes measurable it is predominantly attractive, and both its prevalence and
its strength rise with severity. This is the one finding here that is essentially unchanged
between the two severity scores, because a threshold rather than a gradient drives it.

\begin{figure}[H]
  \centering
  \includegraphics[width=\textwidth]{summary_downloads/Duodenum_clinical_Endocrine_to_Epithelial_summary.pdf}
  \caption{Values follow Figure~\ref{fig:summary}. ``n/a'' in the bottom panel is not a missing
  value: with no eligible control field the rule cannot be evaluated there, so no
  Control-to-disease test exists.}
\end{figure}

\subsection*{Plasma cells lose their place in the colonic immune network}

The loss of IgA-secreting plasma cells is the paper's most emphatic immune finding, and it
matters because those cells maintain the antibody barrier at the gut surface. The rule layer
adds something the counts cannot: attraction falls from 56\% of eligible fields to 0\% while
the plasma share of the tissue falls only from 8.5\% to 3\%, and it happens at Mild rather than
progressively (C--S FDR 0.029).

\paperquote{We found that patients undergoing aHSCT were characterized by a marked decrease in
plasma and CD4 T cells and an increase in various stromal cells, mainly fibroblasts, in
agreement with previous observations.}

\begin{figure}[H]
  \centering
  \includegraphics[width=\textwidth]{summary_downloads/Colon_clinical_CD4T_to_Plasma_summary.pdf}
  \caption{Values follow Figure~\ref{fig:summary}. The Mild group holds 3 eligible fields, which
  the top two panels state under the axis.}
  \label{fig:cd4plasma}
\end{figure}

CD8 T cells do the same thing at the same point, 40\% to 0\%, also at C--S FDR 0.029.

\begin{figure}[H]
  \centering
  \includegraphics[width=\textwidth]{summary_downloads/Colon_clinical_CD8T_to_Plasma_summary.pdf}
  \caption{Values follow Figure~\ref{fig:summary}.}
\end{figure}

This is where the clinical score earns its place as the default: both of these pass under
clinical staging and neither is selected under pathological. The paper reaches its immune
conclusions from the clinical score too.

\subsection*{Stroma takes over as the epithelium retreats}

The paper reports fibrosis and rising fibroblasts alongside epithelial loss. The rule table
shows the stromal side appearing where it was absent: Fibroblast $\rightarrow$ Macrophage is
found in 1 of 56 control fields and in 27 of 115 Mild and 16 of 91 Severe ones
(Figure~\ref{fig:heat}). The corrected duodenal test picks up the matching rearrangement, and
this one holds under both scores: avoidance rises from 6\% to 61\% of eligible fields (ordered
FDR 0.022) while fibroblast and muscle shares both grow, and while the unjudgeable block
\emph{shrinks} from 42\% to 32\% -- so it cannot be an eligibility artefact.

\paperquote{a quantification of cell types across patients undergoing aHSCT showed a high
degree of variability in epithelial cells, with many patients exhibiting a major decrease,
which, in the duodenum, was also accompanied by an increase in fibroblasts and extracellular
matrix (ECM) and correlated with high-stage disease.}

\begin{figure}[H]
  \centering
  \includegraphics[width=\textwidth]{summary_downloads/Duodenum_clinical_Fibroblast_to_Muscle_summary.pdf}
  \caption{Values follow Figure~\ref{fig:summary}.}
  \label{fig:fibmuscle}
\end{figure}

In tissue this is a dense muscle mass with fibroblasts occupying the space around it; the two
barely mix, and the separation deepens as the muscle mass grows.

\begin{figure}[H]
  \centering
  \includegraphics[width=\textwidth]{summary_downloads/Duodenum_clinical_Fibroblast_to_Muscle_fovs.pdf}
  \caption{\fovvalues}
  \label{fig:fovs}
\end{figure}

\subsection*{Colonic goblet cells, and why the score matters}

\hypertarget{scores}{In} healthy colon goblet cells are packed with mucus and hold everything
else out; the screen sees near-total avoidance, with net scores close to $-100\%$. The paper
reports that this population is damaged in colonic aGVHD, and if goblet cells lose their mucus
the exclusion should weaken. It does -- but how it weakens depends entirely on which score is
used, and that is worth stating plainly rather than choosing the friendlier version.

\paperquote{In a subset of colon samples, we observed a loss of mucin expression in the
epithelium, indicating potential damage to goblet cells or impaired mucin production.}

Under pathological staging the relaxation is a clean gradient: avoidance 95\%, 72\%, 27\% of
eligible fields, ordered FDR 0.016. Note the grey unjudgeable block growing from 20\% to 66\%
underneath it.

\begin{figure}[H]
  \centering
  \includegraphics[width=\textwidth]{summary_downloads/Colon_pathological_Goblet_to_Muscle_summary.pdf}
  \caption{Values follow Figure~\ref{fig:summary}; pathological stages.}
  \label{fig:goblet-path}
\end{figure}

Under clinical staging the same rule drops from 95\% to 50\% at Mild and then stops, at 57\% in
Severe. Ordered FDR 0.175, C--M FDR 0.041: a step at the onset of disease rather than a
gradient with severity.

\begin{figure}[H]
  \centering
  \includegraphics[width=\textwidth]{summary_downloads/Colon_clinical_Goblet_to_Muscle_summary.pdf}
  \caption{Values follow Figure~\ref{fig:summary}; the same rule under clinical stages.}
\end{figure}

The reverse direction, calculated independently, moves the same way and by a similar amount:
avoidance 100\%, 67\%, 65\% of eligible fields, C--M FDR 0.029.

\begin{figure}[H]
  \centering
  \includegraphics[width=\textwidth]{summary_downloads/Colon_clinical_Muscle_to_Goblet_summary.pdf}
  \caption{Values follow Figure~\ref{fig:summary}; clinical stages.}
\end{figure}

\subsection*{CD8 T cells pull away from the epithelium}

CD8 T cells are the classical effectors of aGVHD, and in healthy duodenum they are the
population that sits \emph{inside} the epithelium. The paper reports that their distribution
around individual crypts becomes highly uneven in disease. Under pathological staging avoidance
jumps to 75\% of eligible fields in Severe (C--S FDR 0.013) while both populations stay
present -- but Severe is 8 eligible fields, of which 6 avoid.

\paperquote{CD8 T and double-negative T cells (CD3 T cells and double negative for CD4 and
CD8) were the predominant populations inside the epithelium.}

\begin{figure}[H]
  \centering
  \includegraphics[width=\textwidth]{summary_downloads/Duodenum_pathological_CD8T_to_Goblet_summary.pdf}
  \caption{Values follow Figure~\ref{fig:summary}; pathological stages.}
  \label{fig:cd8-path}
\end{figure}

Clinical staging gives Severe 17 eligible fields instead of 8, and the direction repeats --
avoidance 21\%, 22\%, 53\% -- but C--S FDR is 0.315 and nothing passes. The effect is real in
direction and not established in size.

\begin{figure}[H]
  \centering
  \includegraphics[width=\textwidth]{summary_downloads/Duodenum_clinical_CD8T_to_Goblet_summary.pdf}
  \caption{Values follow Figure~\ref{fig:summary}; clinical stages.}
\end{figure}

\subsection*{The plasma-cell arrangement drops at once and never returns}

The paper's temporal result is specific: lymphocytes fall sharply in the first month, myeloid
cells do not, and plasma cells alone fail to recover. Grouping the same colonic fields by time
since transplant reproduces that shape at the level of arrangement -- CD4T $\rightarrow$ Plasma
and CD8T $\rightarrow$ Plasma fall from about $+55$ and $+40$ to zero within 30 days and stay
there past 100 days. Nothing here passes correction, so this is agreement in direction with an
independently established result, not independent evidence.

\paperquote{In contrast, plasma cells, which also dropped shortly after transplantation, failed
to return to normal amounts months after transplantation.}

\begin{figure}[H]
  \centering
  \includegraphics[width=0.97\textwidth]{summary_downloads/time_Colon_paper_immune.pdf}
  \caption{Each small dot is one biopsy and each diamond the group mean; the value is the share
  of that biopsy's eligible FOVs holding attraction minus the share holding avoidance. Control
  sits on the left for context and is not part of the test. The eligible biopsy count is printed
  under each window. These four rules were chosen because the paper highlights these immune
  populations, not because they ranked highly in a screen.}
  \label{fig:time}
\end{figure}

The same four rules stay flat in the duodenum, at temporal FDR 0.99. Both organs are shown so
the contrast is visible rather than selected.

\begin{figure}[H]
  \centering
  \includegraphics[width=0.97\textwidth]{summary_downloads/time_Duodenum_paper_immune.pdf}
  \caption{Values follow Figure~\ref{fig:time}; Duodenum.}
\end{figure}

% =============================================================================
\section{What is new}

None of the results below is stated by the atlas paper. They are ordered by how well they
survive scrutiny, strongest first, because the honest headline of this analysis is how few
results survive both severity scores.

\subsection*{Two rules survive everything}

Of 525 pair rules, exactly two pass the corrected ordered trend test with eligibility applied
under \emph{both} severity scores. Both are duodenal.

\medskip
\noindent\colorbox{soft}{\parbox{0.96\textwidth}{
\textbf{CD4T $\rightarrow$ Goblet} \quad clinical $+28$ pp per stage (FDR 0.022) \quad
pathological $+29$ pp per stage (FDR 0.010)\\
\textbf{Fibroblast $\rightarrow$ Muscle} \quad clinical $-29$ pp per stage (FDR 0.022) \quad
pathological $-28$ pp per stage (FDR 0.010)
}}
\medskip

CD4 T cells and goblet cells start almost completely separated in healthy duodenum and end
barely separated at all: avoidance 77\%, 53\%, 19\% of eligible fields. Why this is
interesting: the paper shows CD4 T cells are zonated along the crypt-villus axis in healthy gut
and fall in number in disease. Here their number is not what changes most -- their access to
the goblet compartment is. CD4T falls from 8.5\% to 3\% of cells while the rule moves three
times as far, and the denominators are the healthiest in the document at 30, 75 and 21 fields.

\begin{figure}[H]
  \centering
  \includegraphics[width=\textwidth]{summary_downloads/Duodenum_clinical_CD4T_to_Goblet_summary.pdf}
  \caption{Values follow Figure~\ref{fig:summary}.}
\end{figure}

In tissue both populations are sparse and interleaved, and the change between stages is
statistical rather than visually dramatic -- which is what a shift of 0.09 in Lift should look
like.

\begin{figure}[H]
  \centering
  \includegraphics[width=\textwidth]{summary_downloads/Duodenum_clinical_CD4T_to_Goblet_fovs.pdf}
  \caption{Selection and rows follow Figure~\ref{fig:fovs}.}
\end{figure}

Fibroblast $\rightarrow$ Muscle was shown in Part 2 (Figures~\ref{fig:fibmuscle}
and~\ref{fig:fovs}) because it also matches the paper's fibrosis result. It is worth restating
why it is unusual: fibroblasts and muscle both grow substantially -- 8\% to 14\% and 4\% to
20\% of cells -- and separate anyway. Abundance and arrangement move in opposite directions,
which is the cleanest demonstration available that the two are different measurements.

\subsection*{Immune and stromal cells pull away from duodenal muscle}

Under pathological staging the muscularis is where the duodenal signal concentrates.
Macrophage $\rightarrow$ Muscle goes from 22\% to 15\% to 70\% avoidance (ordered FDR 0.010,
C--S 0.038) on 18, 66 and 27 eligible fields -- healthy denominators for this cohort -- while
both cell shares roughly double and the unjudgeable block shrinks from 42\% to 25\%. The
Control-to-Severe gap is what carries it; the dip at Mild means the path there is not a clean
gradient.

Why this is interesting: the muscularis has its own resident macrophage population that sits
close to the muscle and to enteric neurons and helps regulate gut motility. That is established
gut biology, independent of transplantation. Losing the macrophage--muscle association is
therefore a candidate readout for muscularis disruption, and it sits alongside the paper's
report of increased fibrosis in the same tissue. What the rules cannot do is identify the
macrophage subtype, so this is a lead about a compartment, not about a cell state.

\begin{figure}[H]
  \centering
  \includegraphics[width=\textwidth]{summary_downloads/Duodenum_pathological_Macrophage_to_Muscle_summary.pdf}
  \caption{Values follow Figure~\ref{fig:summary}; pathological stages.}
\end{figure}

\subsection*{The Paneth niche reorganises, on very few fields}

Read as a whole rather than rule by rule, the pathological duodenal screen concentrates on one
compartment: six of the rules it selects name Paneth cells -- Paneth $\rightarrow$ Mast,
Paneth $\rightarrow$ Macrophage, Paneth $\rightarrow$ Plasma, Paneth $\rightarrow$
Plasma\_CD38, CD8T $\rightarrow$ Paneth and Epithelial $\rightarrow$ Paneth. Paneth cells are
well under 1\% of duodenal cells, so nothing about their abundance would put them at the top of
a screen.

Why this is interesting: Paneth cells sit at the base of the crypt beside the intestinal stem
cells and are the antimicrobial anchor of that niche. The paper proposes that stem cells in
aGVHD may misdifferentiate into endocrine cells instead of Paneth cells. If the crypt-base
niche is being rebuilt, the contacts that define it should change -- and those contacts are
what these rules measure. The rule layer is pointing at the compartment the paper singles out,
through the neighbours rather than the counts.

The caveat comes first and is severe: these rules rest on 3 to 6 eligible Severe fields,
because Paneth cells are rare and become rarer, and none of them appears at all under clinical
staging. Paneth $\rightarrow$ Mast reaches ordered FDR 0.010 on 9, 38 and 3 eligible fields --
so its 67\% Severe avoidance is 2 fields out of 3. This is a lead, not a result.

\begin{figure}[H]
  \centering
  \includegraphics[width=\textwidth]{summary_downloads/Duodenum_pathological_Paneth_to_Mast_summary.pdf}
  \caption{Values follow Figure~\ref{fig:summary}; pathological stages.}
  \label{fig:paneth-sparse}
\end{figure}

Paneth $\rightarrow$ Macrophage adds something the others do not: avoidance both spreads
(8\%, 21\%, 33\%) and deepens (Lift 1.2 to 0.58), at ordered FDR 0.046 on 13, 63 and 6 fields.

\begin{figure}[H]
  \centering
  \includegraphics[width=\textwidth]{summary_downloads/Duodenum_pathological_Paneth_to_Macrophage_summary.pdf}
  \caption{Values follow Figure~\ref{fig:paneth-sparse}.}
\end{figure}

CD8 T cells are the effector population of the disease, so their relation to the crypt base is
the most mechanistically interesting member of the family and the one most worth re-testing in
a larger cohort -- on 12, 38 and 4 eligible fields, M--S FDR 0.025.

\begin{figure}[H]
  \centering
  \includegraphics[width=\textwidth]{summary_downloads/Duodenum_pathological_CD8T_to_Paneth_summary.pdf}
  \caption{Values follow Figure~\ref{fig:paneth-sparse}.}
\end{figure}

\subsection*{Plasma cells and blood vessels swap places in the colon}

Colon Plasma $\rightarrow$ SMV inverts: net $-10\%$, $+25\%$, $+24\%$, with the median Lift
crossing from 0.69 to 1.35 (C--M FDR 0.029). Why this is interesting: plasma cells normally sit
in the lamina propria above the crypts, and a switch from avoiding small vessels to associating
with them is what a population arriving from the circulation rather than residing in the tissue
would look like -- the transition the paper's donor-host chimerism work describes. The Mild
group holds 4 eligible fields, so the middle of the trajectory is thin.

\begin{figure}[H]
  \centering
  \includegraphics[width=\textwidth]{summary_downloads/Colon_clinical_Plasma_to_SMV_summary.pdf}
  \caption{Values follow Figure~\ref{fig:summary}; clinical stages.}
\end{figure}

\subsection*{How strong a rule is, versus how often it is found}

Everything so far counted how often a rule fires. That misses a whole class of change: a rule
found just as often but much stronger or weaker where it is found. Prevalence and Lift are
therefore screened separately, each ranked by its largest stage gap, with no rule named in
advance. The two screens do not return the same rules, which is the point.

Colon Plasma $\rightarrow$ Goblet is the clearest case. Its avoidance Lift drops from 0.73 to
0.22 -- the two types end up roughly three times further apart than chance would place them --
at FDR 0.008, while its prevalence hardly moves at all, at FDR 0.68. A prevalence-only analysis
would have recorded this rule as unchanged.

\begin{figure}[H]
  \centering
  \includegraphics[width=0.88\textwidth]{summary_downloads/Colon_pathological_Plasma_to_Goblet_avoidance_lift.pdf}
  \caption{\statelift{} Colon, pathological stages, avoidance occurrences only. The Severe
  median rests on 3 rule-bearing fields, so the size of the drop is less certain than its
  direction.}
  \label{fig:lift}
\end{figure}

Scored by prevalence instead, the very same rule looks unchanged.

\begin{figure}[H]
  \centering
  \includegraphics[width=0.88\textwidth]{summary_downloads/Colon_pathological_Plasma_to_Goblet_avoidance_prevalence.pdf}
  \caption{The same rule scored by prevalence instead of strength, drawn as the bottom panel of
  Figure~\ref{fig:summary}. Read against Figure~\ref{fig:lift}: two different corrected screens
  reaching opposite conclusions about one rule.}
\end{figure}

The attraction side has a matching case with a clear biological reading. Colon SMV
$\rightarrow$ Endothelial keeps firing in most eligible fields -- net $+73\%$, $+33\%$, $+74\%$,
ordered prevalence FDR 0.978 -- but weakens, median Lift 3.45 to 2.37 at C--S FDR 0.065, and reaches ordered FDR 0.015 under
pathological staging.
Why this is interesting: damage to the endothelium is a recognised feature of transplant
complications, and vessels are one of the few compartments where a loosening of structure has
an obvious physical meaning. This rule measures the association between vessel-lining cells and
the small and medium vessels they belong to, so a weakening association is a plausible spatial
signature of vascular disruption. The atlas does not report it.

\begin{figure}[H]
  \centering
  \includegraphics[width=0.88\textwidth]{summary_downloads/Colon_clinical_SMV_to_Endothelial_attraction_lift.pdf}
  \caption{Values follow Figure~\ref{fig:lift}, for attraction occurrences, where higher Lift
  means stronger association. Clinical stages.}
\end{figure}

Its prevalence, by contrast, is as high in Severe disease as in Control.

\begin{figure}[H]
  \centering
  \includegraphics[width=0.88\textwidth]{summary_downloads/Colon_clinical_SMV_to_Endothelial_attraction_prevalence.pdf}
  \caption{Prevalence for the same rule, drawn as the bottom panel of
  Figure~\ref{fig:summary}.}
\end{figure}



\subsection*{Time after transplant: an honest negative}

The paper makes a strong case that time since transplantation must be modelled, so it is worth
stating what happens when the rules are grouped that way instead of by severity: at biopsy
level, nothing survives correction. In both organs zero rules pass FDR 0.05 in the three-window
screen, and zero pass in either pooled contrast.

The reason is visible in the sampling. The three windows are unevenly filled and the
``more than 100 days'' group spans everything from 100 to over 1100 days, so it is not a
homogeneous group at all.

\begin{figure}[H]
  \centering
  \includegraphics[width=0.85\textwidth]{summary_downloads/time_coverage.pdf}
  \caption{Each dot is one biopsy at the number of days between transplantation and biopsy,
  coloured by the window it was assigned to. Biopsies are the unit of every temporal test, so
  this is the effective sample size of this whole subsection.}
\end{figure}

Screening every rule against the three windows returns nothing at all in the colon.

\begin{figure}[H]
  \centering
  \includegraphics[width=0.90\textwidth]{summary_downloads/time_Colon_screen.pdf}
  \caption{Each dot is one rule: the horizontal axis is the largest gap between any two
  post-transplant windows, the vertical axis the FDR across all rules tested at biopsy level,
  drawn so stronger results sit higher. Colon.}
  \label{fig:temporal}
\end{figure}

The largest duodenal signals are all plasma-cell and goblet-cell rules, which is at least the
compartment the paper's kinetics point to, but none of them passes either.

\begin{figure}[H]
  \centering
  \includegraphics[width=0.90\textwidth]{summary_downloads/time_Duodenum_screen.pdf}
  \caption{Values follow Figure~\ref{fig:temporal}; Duodenum.}
\end{figure}

Pooling the windows into two broader comparisons does not rescue it.

\begin{figure}[H]
  \centering
  \includegraphics[width=0.90\textwidth]{summary_downloads/time_Colon_pooled.pdf}
  \caption{Two contrasts, each read as in Figure~\ref{fig:temporal}, with one FDR correction
  applied jointly across both contrasts and all their rules.}
\end{figure}

% =============================================================================
\section{What to take from this}

\textbf{The method measures arrangement, not composition.} Across every comparison in Part 1
the rules move several times further than the cell counts behind them -- 67 percentage points
against 9 in the clearest case -- and Fibroblast $\rightarrow$ Muscle separates while both of
its cell types grow.

\textbf{Eligibility control is not a detail, it is most of the analysis.} Without it, 22
colonic rules climb and 11 fade with clinical severity. With it, none of them passes. A rule
going from avoidance to unjudgeable looks exactly like avoidance relaxing, and in the colon --
where goblet cells fall from 22\% to 13\% of the tissue -- that is what most of the apparent
signal was.

\textbf{Where the atlas has a finding, the rule layer recovers it independently:} Paneth loss,
endocrine accumulation, the collapse of the colonic plasma-cell network, the stromal shift, the
organ split, and the immediate and permanent loss of plasma-cell organisation after transplant.

\textbf{Two results survive everything.} Duodenal CD4T $\rightarrow$ Goblet and Fibroblast
$\rightarrow$ Muscle pass the corrected ordered trend test with eligibility applied under both
severity scores, with the best denominators in the cohort. Everything else in Part 3 is
score-dependent: the Paneth family and the monotonic colonic goblet gradient are pathological
findings, the colonic plasma attractions are clinical findings.

\textbf{That score-dependence is itself the most interesting thing here.} It is not noise. The
epithelial results track the pathological score, which is read off the same slide the rules are
measured in; the immune results track the clinical score, which is the one the paper uses for
its own immune work. The paper notes that the two gradings are only moderately correlated and
suggests imaging features might bridge them. These rules behave exactly as though the two
scores are measuring different things.

\textbf{The limits.} Many endpoints rest on a handful of eligible fields -- the Paneth family
worst, at 3 to 6. The temporal analysis is a clean negative at biopsy level. And every result
here is a spatial association between classifier labels, which is where a mechanism starts
rather than evidence of one.

\end{document}
